\section{Capitolo 4: TypeScript e Tipi}

\begin{esercizio}{Es. 13: Interfacce e Tipi Opzionali}
\textbf{TRACCIA:} Definire un'interfaccia `Studente` con:
- `matricola` (numero, obbligatorio)
- `nome` (stringa, obbligatorio)
- `votoLaurea` (numero, \textbf{opzionale} - potrebbe non essersi ancora laureato)

Scrivere poi una funzione `isLaureato(s)` che ritorna true se il voto esiste.
\end{esercizio}

\begin{jsbox}
interface Studente {
    matricola: number;
    nome: string;
    votoLaurea?: number; // Il punto interrogativo indica opzionale
}

function isLaureato(s: Studente): boolean {
    // Controllo se è diverso da undefined
    return s.votoLaurea !== undefined;
}
\end{jsbox}

\begin{esercizio}{Es. 14: Generics (La Scatola Magica)}
\textbf{TRACCIA:} Scrivere una classe `Scatola<T>` che può contenere un valore di qualsiasi tipo T.
Metodi: `inserisci(val: T)` e `estrai(): T`.
\end{esercizio}

\begin{jsbox}
class Scatola<T> {
    private contenuto: T | null = null;
    
    inserisci(val: T): void {
        this.contenuto = val;
    }
    
    estrai(): T | null {
        return this.contenuto;
    }
}

// Uso
let scatolaNumeri = new Scatola<number>();
scatolaNumeri.inserisci(10);
// scatolaNumeri.inserisci("ciao"); // ERRORE TS!
\end{jsbox}

\begin{esercizio}{Es. 15: Tuple e Array}
\textbf{TRACCIA:} Definire un tipo `Punto3D` come una Tripla (Tupla) di 3 numeri `[x, y, z]`.
Scrivere una funzione che somma due punti.
\end{esercizio}

\begin{jsbox}
type Punto3D = [number, number, number];

function sommaPunti(p1: Punto3D, p2: Punto3D): Punto3D {
    return [
        p1[0] + p2[0],
        p1[1] + p2[1],
        p1[2] + p2[2]
    ];
}
\end{jsbox}

\begin{esercizio}{Es. 16: Type Guards (Narrowing)}
\textbf{TRACCIA:} Una funzione accetta `input: string | number`.
Se è una stringa, ritorna la lunghezza. Se è un numero, ritorna il numero stesso moltiplicato per 2.
Usa `typeof` per distinguere.
\end{esercizio}

\begin{jsbox}
function processa(input: string | number): number {
    if (typeof input === "string") {
        // Qui TS sa che input è string
        return input.length;
    } else {
        // Qui TS sa che input e' number
        return input * 2;
    }
}
\begin{esercizio}{Es. 25: Discriminated Union (Result Type)}
\textbf{TRACCIA:} Definire un tipo `Result<T>` che può essere `{ success: true, data: T }` OPPURE `{ success: false, error: string }`.
Scrivere una funzione che gestisce entrambi i casi stampando il dato o l'errore.
\end{esercizio}

\begin{jsbox}
type Result<T> = 
    | { success: true; data: T } 
    | { success: false; error: string };

function gestisci(res: Result<number>) {
    if (res.success) {
        // Qui TS sa che esiste .data
        console.log("Dato:", res.data);
    } else {
        // Qui TS sa che esiste .error
        console.log("Errore:", res.error);
    }
}
\end{jsbox}
